\chapter{JADWAL PENELITIAN}

Penelitian ini direncanakan untuk dilaksanakan selama enam bulan, dari bulan Februari 2026 hingga Juli 2026. Jadwal penelitian disusun berdasarkan tahapan metodologi Research and Development (R\&D) yang telah dijelaskan pada Bab III.

Tabel \ref{tab:jadwal} menunjukkan rencana jadwal pelaksanaan penelitian berbasis bulan. Simbol $\bullet$ menunjukkan bahwa kegiatan dilaksanakan pada bulan tersebut.

\begin{table}[H]
  \centering
  \caption{Jadwal Penelitian}
  \label{tab:jadwal}
  \begin{tabular}{|l|c|c|c|c|c|c|}
    \hline
    \textbf{Kegiatan} & \textbf{Feb} & \textbf{Mar} & \textbf{Apr} & \textbf{Mei} & \textbf{Jun} & \textbf{Jul} \\
    \hline
    Studi Literatur & $\bullet$ & $\bullet$ & & & & \\
    \hline
    Analisis Kebutuhan & $\bullet$ & $\bullet$ & & & & \\
    \hline
    Perancangan Sistem & & $\bullet$ & $\bullet$ & & & \\
    \hline
    Implementasi Sistem & & & $\bullet$ & $\bullet$ & & \\
    \hline
    Pengujian dan Evaluasi & & & & $\bullet$ & $\bullet$ & \\
    \hline
    Dokumentasi dan Penulisan & & $\bullet$ & $\bullet$ & $\bullet$ & $\bullet$ & $\bullet$ \\
    \hline
    Seminar Proposal / Sidang & & & & & & $\bullet$ \\
    \hline
  \end{tabular}
\end{table}

\section{Penjelasan Kegiatan}

\begin{enumerate}
    \item \textbf{Studi Literatur (Februari -- Maret 2026)}
    
    Kegiatan ini mencakup pengumpulan dan analisis literatur mengenai InvenioRDM, Kubernetes, GitOps, ArgoCD, dan pola Apps of Apps. Sumber literatur meliputi dokumentasi resmi, buku teks, jurnal ilmiah, dan artikel teknis.
    
    \item \textbf{Analisis Kebutuhan (Februari -- Maret 2026)}
    
    Tahap ini meliputi analisis kebutuhan komponen InvenioRDM, dependensi antar komponen, kebutuhan resource Kubernetes, serta identifikasi tools dan teknologi yang diperlukan untuk implementasi.
    
    \item \textbf{Perancangan Sistem (Maret -- April 2026)}
    
    Kegiatan perancangan meliputi pembuatan diagram arsitektur sistem, desain struktur repositori Git, perancangan hierarki Apps of Apps ArgoCD, dan perancangan manifest Kubernetes serta Helm chart.
    
    \item \textbf{Implementasi Sistem (April -- Mei 2026)}
    
    Tahap implementasi meliputi pembuatan klaster Kubernetes, instalasi ArgoCD, pembuatan repositori Git, implementasi aplikasi ArgoCD induk dan anak, konfigurasi Helm chart InvenioRDM, serta konfigurasi networking dan storage.
    
    \item \textbf{Pengujian dan Evaluasi (Mei -- Juni 2026)}
    
    Kegiatan pengujian meliputi pengujian fungsional deployment, pengujian sinkronisasi GitOps, pengujian drift detection, pengujian rollback, serta pengukuran metrik kinerja sistem.
    
    \item \textbf{Dokumentasi dan Penulisan (Maret -- Juli 2026)}
    
    Kegiatan dokumentasi dilakukan secara paralel dengan tahapan lainnya. Penulisan proposal tugas akhir dilakukan pada semester ini, sedangkan penulisan skripsi lengkap akan dilanjutkan pada semester berikutnya.
    
    \item \textbf{Seminar Proposal / Sidang (Juli 2026)}
    
    Tahap akhir adalah presentasi dan sidang proposal tugas akhir di depan dosen penguji.
\end{enumerate}
